\documentclass[a4paper,12pt]{article}

% Språk og tegnsett
\usepackage{fontspec}
\setmainfont{Times New Roman}
\usepackage{xcolor}  

\usepackage[norsk]{babel}
\usepackage{csquotes}

% Marger og layout
\usepackage[a4paper,margin=2.5cm]{geometry}

% APA-sitering via biblatex
\usepackage[style=apa, backend=biber]{biblatex}

% Last inn referansefil 
\addbibresource{referanser.bib}

\usepackage{setspace}
\onehalfspacing

% For lenker og DOI
%\usepackage[hidelinks]{hyperref}
\usepackage{hyperref}

\hypersetup{
    colorlinks=true,
    linkcolor=black,        % Seksjonslenker og TOC
    citecolor=red!50!black, % Referanser
    urlcolor=blue           % Nettadresser
}

\usepackage{tocloft}

% Normal skrift i TOC
\renewcommand{\cftsecfont}{\normalfont}
\renewcommand{\cftsubsecfont}{\normalfont}
\renewcommand{\cftsecpagefont}{\normalfont}
\renewcommand{\cftsubsecpagefont}{\normalfont}

\DefineBibliographyStrings{norsk}{
  andothers = {et al.},
  editor    = {red.},
  editors   = {red.},
  in        = {i},
  edition   = {utg.},
  urlseen   = {Hentet fra},
  and = {og},
  %page = {s.},
  %pages = {ss.},
}

% For norsk dato
%\usepackage{datetime2}
%\DTMsetup{locale=nb-NO}

%\usepackage{titlesec}

% Seksjoner – litt større enn små
%\titleformat{\section}
%  {\normalfont\large\bfseries}{\thesection}{1em}{}

% Underseksjoner – behold små
%\titleformat{\subsection}
%  {\normalfont\normalsize\bfseries}{\thesubsection}{1em}{}

\title{Analyse av undervisningsopplegg}
\author{Christian Lo Virik}
\date{\today}

\begin{document}

\maketitle

\abstract{
\noindent 
 I denne oppgaven analyseres et undervisningsopplegg i naturfag fra førsteklasse studieforberedende på videregående skole. Undervisningsopplegget analyseres med hensyn til prinsippet om tilpasset opplæring og flerkulturell pedagogikk. Selv om opplegget viser gode sider .. ville noen enkle justeringer gjort det mye bedre.
} 	
\newpage
\tableofcontents
\newpage

\hypersetup{
    colorlinks=true,
    linkcolor=red!50!black,        % Seksjonslenker og TOC
    citecolor=red!50!black, % Referanser
    urlcolor=blue           % Nettadresser
}





\section{Introduksjon}
Denne oppgaven tar utgangspunkt i et konkret undervisningsopplegg i naturfag på Vg1 med tema ernæring og helse, gjennomført i en dobbelttime med en sammensatt elevgruppe. Formålet er å analysere og drøfte opplegget i lys av relevant pedagogisk teori, samt gjeldende lovverk og forskrifter, med særlig vekt på prinsipper om tilpasset opplæring og inkludering. Klassen bestod av elever med ulike forutsetninger, blant annet minoritetsspråklige elever, elever med lese- og skrivevansker, elever med individuell opplæringsplan (IOP) og elever med høyt faglig nivå. Dette gjør undervisningssituasjonen godt egnet for å undersøke hvordan læreren kan legge til rette for læring i et mangfoldig klasserom.

Oppgaven avgrenses til å fokusere på hvordan tilpasset opplæring kommer til uttrykk i planleggingen og gjennomføringen av undervisningsopplegget, særlig gjennom bruk av stasjonsundervisning og varierte arbeidsmåter. I tillegg vil oppgaven belyse hvordan tiltakene retter seg mot ulike elevgrupper, og i hvilken grad disse kan bidra til inkludering og faglig utvikling. Perspektiver fra flerkulturell pedagogikk vil også trekkes inn for å vurdere hvordan undervisningen ivaretar minoritetsspråklige elevers behov for språkstøtte og deltakelse.

Videre vil analysen knyttes til sentrale bestemmelser i opplæringsloven og prinsipper i læreplanverket, som understreker skolens ansvar for å gi alle elever et likeverdig opplæringstilbud. Relevant teori om differensiering, tilpasset opplæring og inkluderende praksis vil danne grunnlag for drøftingen. Det vil ikke bli gjort en fullstendig vurdering av alle sider ved undervisningen, men hovedvekten legges på de didaktiske valgene som omhandler tilpasning og inkludering.

Gjennom denne avgrensningen søker oppgaven å gi en faglig fundert vurdering av hvordan undervisningsopplegget kan forstås, forbedres og videreutvikles i lys av pedagogisk teori og praksis.

\section{Undervisningsopplegget}
Undervisningsopplegget som analyseres i denne oppgaven ble gjennomført i løpet av en dobbelttime i naturfag på Vg1, temaet for timen var ernæring og helse. I utviklingen av opplegget ble det lagt vekt på tilpasset opplæring, siden klassen bestod av en variert elevgruppe med ulike faglige, språklige og sosiale utfordringer. Det ble tatt hensyn til minoritetsspråklige elever, elever med lese- og skrivevansker, elever med IOP og også høyt presterende elever.

I starten av timen ble det lagt til rette for trygge og forutsigbare rammer ved felles gjennomgang av samspillsregler. Etter på ble elevenes forkunnskaper aktivert ved en introduksjonsvideo etterfulgt av en lærerstyrt samtale. De minoritetsspråklige elevene fikk støtte gjennom bildeordlister som økte deltakelsen i faglige samtaler.

Hoveddelen av undervisningen var organisert som stasjonsundervisning med fire læringsaktiviteter. Dette ga mulighet for tilpasning både i nivå og støtte.

På første stasjon skulle elevene lese en kort fagtekst og svar på fem refleksjonsspørsmål. Her kunne elevene fritt velge en lettlest variant om de hadde behov for det. De minoritetsspråklige elevene jobbet i par slik at de kunne få hjelp fra andre elver med samme bakgrunn. Elevene med lese- og skrivevansker fikk lov til å bruke opplesningsfunksjon og også svare muntlig på refleksjonsoppgavene. Elevene med IOP fikk forenklet tekst med nøkkelsetninger markert for å gjøre innholdet enklere å forstå.

Andre stasjon handlet om aminosyrenes struktur, og elevene skulle bygge aminosyrer med et molekylbyggesett. Elevene med IOP brukte et forenklet byggesett, mens de faglig sterke elevene fikk noen fordypningsoppgaver om effekten av rekkefølgen av aminosyrene på proteinets struktur på dets funksjon.

Tredje stasjon fokuserte på visuelle framstillinger av proteiner og hvor i kroppen de finnes og hva de brukes til. Blideordlister ble lagt frem for å hjelpe de minoritetsspråklige elevene og elevene med lese- og skrivevansker. Disse fikk også lov til å svare muntlig. De faglig sterke elevene fikk lese om proteinenes funksjoner på molekylært nivå.

Fjerde stasjon bestod av et praktisk forsøk om denaturering av proteiner. Her skulle elevene fylle ut et observasjonsskjema.  Eleven med IOP brukte et visuelt avkrysningsskjema. De faglig sterke elevene fikk også lese en fordypende tekst om denaturering og renaturering.

Til slutt ble det gjennomført en felles oppsummering og elevene fikk en individuell innlevering i lekse. Elevene skulle velge en modell eller representasjon for undervisningen og forklare hva denne viser oss. De kunne velge mellom skriftlig og muntlig innlevering. \parencite[§5]{opplaeringsloven}.

\section{Flerkulturell pedagogikk}
\cite{iversen2025}
\section{Tilpasset opplæring}
\cite{holland2021}
\section{Opplæringsloven}

\section{Konklusjon}

\section*{Merknad om bruk av kunstig intelligens}

\textcite{openai_chatgpt} sier det kanskje best selv: 
\begin{quote}
I arbeidet med denne oppgaven er kunstig intelligens (KI) -- nærmere bestemt ChatGPT, utviklet av OpenAI -- benyttet som et støtteverktøy i forbindelse med rettskrivning, språklig bearbeiding og forbedring av enkelte formuleringer. Alt faglig innhold som ikke er eksplisitt referert til, herunder analyser, vurderinger og konklusjoner, er utelukkende utarbeidet av forfatteren selv.
\end{quote}

%\bibliographystyle{apacite}
\addcontentsline{toc}{section}{Litteraturliste}
\printbibliography[title={Litteraturliste}]
\end{document}